\section{Chamber stabilizers and the splitting issue}
A geometry is \emph{thin} when every panel belongs to exactly two chambers.
The chamber graph joins chambers differing in one object. This graph need not
be connected merely from the definition that every flag extends to a chamber.

\begin{proposition}[The connected case]\label{UR.ThinSplit}
Suppose a thin geometry has a connected chamber graph and its type-preserving
automorphism group $K$ is chamber-transitive. Write $C$ for its full correlation
group. Then $K$ acts regularly on chambers. For any chamber $D$, its setwise
stabilizer $C_D$ is a complement to $K$ in $C$:
\[C=K\rtimes C_D,\qquad C_D\cong C/K.\]
\end{proposition}
\begin{proof}
An element of $K$ fixing $D$ fixes its objects individually. It fixes each
adjacent chamber: the common panel is fixed, and of its two chambers one is
already fixed, so the other is fixed as well. Connectedness propagates this
to every chamber. Every object lies in a chamber; hence the element is the
identity. Chamber transitivity now gives regularity. Given $c\in C$, choose
$k\in K$ with $kD=cD$. Then $k^{-1}c\in C_D$. Thus $C=KC_D$, while
$K\cap C_D=1$. Normality of the palette kernel $K$ proves the semidirect
product and the quotient identification.
\end{proof}

\begin{proposition}[Thin doubling of a bipartite graph]\label{UR.ThinDoubling}
Let $B=(P\sqcup Q,E)$ be a finite connected simple bipartite graph with
$|P|=|Q|\ne|E|$. Assume that its part-preserving automorphism group
$H^+$ is transitive on $E$ and that some automorphism exchanges its parts.
Write $H=\Aut(B)$, so $H/H^+\cong C_2$. There is a finite thin,
flag-transitive rank-three geometry $\Delta(B)$ with
\[
 C(\Delta(B))\cong C_2^{P\sqcup Q\sqcup E}\rtimes H,
 \qquad
 K(\Delta(B))\cong C_2^{P\sqcup Q\sqcup E}\rtimes H^+.
\]
If the exchanging coset $H\setminus H^+$ contains no involution, the
first group does not split over the second. Its chamber graph has one
component for each edge of $B$.
\end{proposition}
\begin{proof}
For each $v\in P\sqcup Q\sqcup E$ take two objects $(v,0),(v,1)$,
with the three displayed sets as their types. Incidence ignores the
second coordinate: a $P$-object and a $Q$-object are incident precisely
when their underlying vertices form an edge of $B$; an edge-object is
incident with precisely the objects over its two endpoints. Distinct objects of one type are nonincident; we use the source's
irreflexive incidence convention.
A chamber is exactly an edge together with independent choices of one
of the two objects over each endpoint and that edge. Every incident
pair extends to such a chamber, as does every single object since the
graph is connected and has an edge. Thus every flag extends to a
chamber. For a panel, the underlying missing vertex or edge is unique,
and its two copies give exactly two completions. This proves thinness.

The independent flips of all pairs form
$N=C_2^{P\sqcup Q\sqcup E}$. The group $H$ lifts by leaving second
coordinates unchanged. The subgroup $N\rtimes H^+$ is transitive on
chambers, using edge transitivity followed by the three endpoint/edge
flips. It is therefore transitive on flags of each type.

Every correlation fixes the edge type: its size $2|E|$ differs from
both other type sizes. It either fixes or exchanges the other two
types. Within each type the two copies of an underlying object are
exactly the classes with equal neighborhoods in the other types.
For distinct vertices of $P$ or $Q$, their sets of incident edge-objects
are disjoint and nonempty, so their neighborhoods differ. Distinct
edge-objects have distinct pairs of endpoints because $B$ is simple.
Consequently these pairs are intrinsically determined, and every
correlation descends to an automorphism of $B$, possibly exchanging
its parts. Its action on edges is forced by its action on vertices.
The restriction kernel consists exactly of the independent flips $N$.
The explicit lifts split this restriction and prove both group formulas.

The quotient by the type-preserving subgroup is $C_2$. A complement
would contain an involution exchanging the two vertex types. Its
projection to $H$ would be an involution in $H\setminus H^+$,
contrary to the stated hypothesis. Finally adjacent chambers have the
same underlying edge; the eight chambers over each edge form the
connected three-dimensional cube under changes of a single choice.
These are precisely the chamber components.
\end{proof}

The second application supplies a finite graph satisfying the nonsplitting
hypothesis, with a complete elementary certificate of its full automorphism
group. Thus the unqualified source observation is false, whereas the
connected-chamber statement above holds by its direct proof.
